\documentclass[10pt,a4paper]{amsart}
\usepackage{amsmath,amssymb,amsthm,graphicx}
% Retained minimal preamble for the Stacks Project chapter sites-cohomology.tex (Cohomology on Sites).
% Verbatim excerpt of the chapter preamble: Xy-pic setup (native preamble lines 14-18) and the native
% theorem declarations with the mathematical macros (native preamble lines 154-217). Copied mechanically;
% nothing added, removed or redefined. The wrapper class and the amsmath/amssymb/amsthm/graphicx base
% packages are supplied by the assembler, so the class/externaldocument part of the native preamble is omitted.
% For commutative diagrams we use Xy-pic
\usepackage[all]{xy}

% We use 2cell for 2-commutative diagrams.
\xyoption{2cell}
% Theorem environments.
%
\theoremstyle{plain}
\newtheorem{theorem}[subsection]{Theorem}
\newtheorem{proposition}[subsection]{Proposition}
\newtheorem{lemma}[subsection]{Lemma}

\theoremstyle{definition}
\newtheorem{definition}[subsection]{Definition}
\newtheorem{example}[subsection]{Example}
\newtheorem{exercise}[subsection]{Exercise}
\newtheorem{situation}[subsection]{Situation}

\theoremstyle{remark}
\newtheorem{remark}[subsection]{Remark}
\newtheorem{remarks}[subsection]{Remarks}

\numberwithin{equation}{subsection}

% Macros
%
\def\lim{\mathop{\mathrm{lim}}\nolimits}
\def\colim{\mathop{\mathrm{colim}}\nolimits}
\def\Spec{\mathop{\mathrm{Spec}}}
\def\Hom{\mathop{\mathrm{Hom}}\nolimits}
\def\Ext{\mathop{\mathrm{Ext}}\nolimits}
\def\SheafHom{\mathop{\mathcal{H}\!\mathit{om}}\nolimits}
\def\SheafExt{\mathop{\mathcal{E}\!\mathit{xt}}\nolimits}
\def\Sch{\mathit{Sch}}
\def\Mor{\mathop{\mathrm{Mor}}\nolimits}
\def\Ob{\mathop{\mathrm{Ob}}\nolimits}
\def\Sh{\mathop{\mathit{Sh}}\nolimits}
\def\NL{\mathop{N\!L}\nolimits}
\def\CH{\mathop{\mathrm{CH}}\nolimits}
\def\proetale{{pro\text{-}\acute{e}tale}}
\def\etale{{\acute{e}tale}}
\def\QCoh{\mathit{QCoh}}
\def\Ker{\mathop{\mathrm{Ker}}}
\def\Im{\mathop{\mathrm{Im}}}
\def\Coker{\mathop{\mathrm{Coker}}}
\def\Coim{\mathop{\mathrm{Coim}}}

% Boxtimes
%
\DeclareMathSymbol{\boxtimes}{\mathbin}{AMSa}{"02}

%
% Macros for moduli stacks/spaces
%
\def\QCohstack{\mathcal{QC}\!\mathit{oh}}
\def\Cohstack{\mathcal{C}\!\mathit{oh}}
\def\Spacesstack{\mathcal{S}\!\mathit{paces}}
\def\Quotfunctor{\mathrm{Quot}}
\def\Hilbfunctor{\mathrm{Hilb}}
\def\Curvesstack{\mathcal{C}\!\mathit{urves}}
\def\Polarizedstack{\mathcal{P}\!\mathit{olarized}}
\def\Complexesstack{\mathcal{C}\!\mathit{omplexes}}
% \Pic is the operator that assigns to X its picard group, usage \Pic(X)
% \Picardstack_{X/B} denotes the Picard stack of X over B
% \Picardfunctor_{X/B} denotes the Picard functor of X over B
\def\Pic{\mathop{\mathrm{Pic}}\nolimits}
\def\Picardstack{\mathcal{P}\!\mathit{ic}}
\def\Picardfunctor{\mathrm{Pic}}
\def\Deformationcategory{\mathcal{D}\!\mathit{ef}}
\usepackage{cleveref}
\crefname{originallemma}{original lemma}{original lemmas}
\Crefname{originallemma}{Original lemma}{Original lemmas}
\makeatletter
\newlabel{lemma-downstairs}{{0D7Q}{source}{}{sourcepointer-1}{}}
\newlabel{lemma-downstairs@cref}{{[originallemma][0][]0D7Q}{[1][0][]source}{}{}{}}
\makeatother
\title{A bounded below object fixed by pullback-pushforward along a flat morphism of ringed topoi}
\author{The Stacks Project Authors}
\date{stacks-project:sites-cohomology.tex:local-snapshot-20261011; GNU FDL 1.2 or later}
\begin{document}
\maketitle
\noindent\textit{Editorial note.} All delivered mathematical text is the author's own, copied from the pinned Stacks Project chapter file sites-cohomology.tex; no statement or proof step was rewritten, reordered or removed.\par\medskip
\noindent\textit{Editorial note.} The proof is the author's complete reduction: the passage to cohomology sheaves, the truncation identity and the membership in the subcategory D' are all in the counted span.\par\medskip
\noindent\textit{Editorial note.} The label lemma-bounded-in-image-downstairs occurs in no other staged Stacks chapter, so no cross-chapter restatement had to be weighed for this card.\par\medskip
\noindent\textit{Editorial note.} The single labelled reference occurring inside the counted spans is carried in assembly.external\_references as a non-load-bearing navigation pointer to the same chapter.\par\medskip
\noindent\textit{Editorial note.} The preamble is the chapter's own retained minimal preamble (sources/stacks-sites-cohomology/retained-preamble.tex), a verbatim excerpt of the native preamble: the Xy-pic setup and the native theorem declarations with the mathematical macros.\par\medskip
\noindent\textit{Allowed standard prerequisites (editorial).}\begin{itemize}
\item Flat morphism of ringed topoi. Same chapter, the definition of a flat morphism of ringed topoi. Supplies the hypothesis under which the truncation computation H\textasciicircum{}j(Rf\_*f\textasciicircum{}*K) = H\textasciicircum{}j(Rf\_*f\textasciicircum{}*tau\_\{<= j\}K) is carried out.
\item Cohomology sheaves and truncations tau\_\{<= j\} of an object of D(O\_D). Same chapter, the truncation functors and the criterion that an object is zero if and only if all its cohomology sheaves vanish. Supplies the reduction to the bounded case and the identity H\textasciicircum{}j(Rf\_*f\textasciicircum{}*K) = H\textasciicircum{}j(Rf\_*tau\_\{<= j\}f\textasciicircum{}*K).
\item The triangulated subcategory of objects whose cohomology sheaves satisfy H\textasciicircum{}q -> Rf\_*f\textasciicircum{}*H\textasciicircum{}q. Same chapter, the subcategory D' introduced in the downstairs lemma and its stability under distinguished triangles. Supplies the final step, in which K is placed in D' because all its cohomology sheaves are.
\end{itemize}\medskip
\ifdefined\hypertarget\hypertarget{sourcepointer-1}{}\fi
\noindent\textit{Source reference.} Original lemma 0D7Q: sources/stacks-sites-cohomology/pinned-original-source.tex, source lines 6286--6286. Same-chapter navigation pointer invoked at proof line 6387, where the author places the cohomology sheaves in the triangulated subcategory introduced by that earlier lemma. The pointer only resolves that reference for the reader; the cited lemma is not reproduced inside the counted span, and no bibliography citation is involved.\par\medskip
\begin{lemma}
\label{lemma-bounded-in-image-downstairs}
Let $f : (\Sh(\mathcal{C}), \mathcal{O}_\mathcal{C}) \to
(\Sh(\mathcal{D}), \mathcal{O}_\mathcal{D})$ be a morphism of ringed topoi.
Let $K$ be an object of $D(\mathcal{O}_\mathcal{D})$. Assume
\begin{enumerate}
\item $f$ is flat,
\item $K$ is bounded below,
\item $H^q(K) \to Rf_*f^*H^q(K)$ is an isomorphism.
\end{enumerate}
Then $K \to Rf_*f^*K$ is an isomorphism.
\end{lemma}

\begin{proof}
Observe that $K \to Rf_*f^*K$ is an isomorphism if and only
if it is an isomorphism on cohomology sheaves $H^j$. Observe that
$H^j(Rf_*f^*K) = H^j(Rf_*\tau_{\leq j}f^*K) = H^j(Rf_*f^*\tau_{\leq j}K)$.
Hence we may assume that $K$ is bounded. Then property (3)
tells us the cohomology sheaves are in the triangulated
subcategory $D' \subset D(\mathcal{O}_\mathcal{D})$ of
Lemma \ref{lemma-downstairs}. Hence $K$ is in it too.
\end{proof}

\end{document}
